\documentclass[aspectratio=169,11pt]{beamer}
% Lecture 01 — First program, console, variables, templates
% Course: Introduction to Programming with Kotlin

% Shared Beamer preamble for Introduction to Programming with Kotlin
% Include with: % Shared Beamer preamble for Introduction to Programming with Kotlin
% Include with: % Shared Beamer preamble for Introduction to Programming with Kotlin
% Include with: \input{preamble}

\usetheme{Madrid}
\usecolortheme{default}

% Clean, professional palette (deep navy + muted teal accent)
\definecolor{LectureNavy}{HTML}{1B3A4B}
\definecolor{LectureAccent}{HTML}{2A9D8F}
\definecolor{LectureGray}{HTML}{4A5568}
\definecolor{CodeBg}{HTML}{F7FAFC}
\definecolor{AlertBg}{HTML}{FFF5F5}
\definecolor{TryBg}{HTML}{F0FFF4}
\definecolor{QuizBg}{HTML}{EBF8FF}
\definecolor{AnswerKeyBg}{HTML}{FFFFF0}
\definecolor{AnswerGold}{HTML}{B7791F}
\definecolor{KeywordBlue}{HTML}{2B6CB0}
\definecolor{StringGreen}{HTML}{276749}
\definecolor{CommentGray}{HTML}{718096}

\setbeamercolor{structure}{fg=LectureNavy}
\setbeamercolor{palette primary}{bg=LectureNavy,fg=white}
\setbeamercolor{palette secondary}{bg=LectureAccent,fg=white}
\setbeamercolor{palette tertiary}{bg=LectureGray,fg=white}
% Madrid puts title/subtitle in a coloured box (inherits palette primary bg).
% Title text must be light on that dark box — not navy-on-navy.
\setbeamercolor{title}{fg=white,bg=LectureNavy}
\setbeamercolor{subtitle}{fg=white,bg=LectureNavy}
\setbeamercolor{author}{fg=LectureNavy}
\setbeamercolor{date}{fg=LectureGray}
\setbeamercolor{institute}{fg=LectureGray}
\setbeamercolor{frametitle}{bg=LectureNavy,fg=white}
\setbeamercolor{block title}{bg=LectureNavy,fg=white}
\setbeamercolor{block body}{bg=CodeBg,fg=black}
\setbeamercolor{block title alerted}{bg=LectureAccent,fg=white}
\setbeamercolor{block body alerted}{bg=TryBg,fg=black}
\setbeamercolor{block title example}{bg=LectureGray,fg=white}
\setbeamercolor{block body example}{bg=AlertBg,fg=black}

\setbeamertemplate{navigation symbols}{}
\setbeamertemplate{itemize items}[circle]
\setbeamertemplate{enumerate items}[default]

\usepackage[T1]{fontenc}
\usepackage[utf8]{inputenc}
\usepackage{lmodern}
\usepackage{booktabs}
\usepackage{tabularx}
\usepackage{graphicx}
\usepackage{hyperref}
\usepackage{listings}
\usepackage{xcolor}
\usepackage{tikz}
\usetikzlibrary{positioning,arrows.meta}

% listings: Kotlin without shell-escape (no built-in Kotlin language)
\lstdefinelanguage{Kotlin}{
  morekeywords={
    abstract,actual,annotation,as,break,by,catch,class,companion,const,
    constructor,continue,crossinline,data,do,dynamic,else,enum,expect,
    external,false,field,file,final,finally,for,fun,get,if,import,in,
    infix,init,inline,inner,interface,internal,is,lateinit,noinline,null,
    object,open,operator,out,override,package,private,protected,public,
    reified,return,sealed,set,super,suspend,tailrec,this,throw,true,try,
    typealias,typeof,val,var,when,where,while
  },
  sensitive=true,
  morecomment=[l]{//},
  morecomment=[s]{/*}{*/},
  morestring=[b]",
  morestring=[b]',
  morestring=[s]{"""}{"""},
}

\lstdefinestyle{kotlinstyle}{
  language=Kotlin,
  basicstyle=\ttfamily\small,
  keywordstyle=\color{KeywordBlue}\bfseries,
  commentstyle=\color{CommentGray}\itshape,
  stringstyle=\color{StringGreen},
  numbers=left,
  numberstyle=\tiny\color{CommentGray},
  numbersep=6pt,
  frame=single,
  rulecolor=\color{LectureGray!40},
  backgroundcolor=\color{CodeBg},
  breaklines=true,
  showstringspaces=false,
  tabsize=2,
  captionpos=b,
  morekeywords={readln,readLine,println,print,listOf,mutableListOf,toInt,toDouble,toLong,toShort,toByte,toChar,toString,split,size},
}
\lstset{style=kotlinstyle}

\newcommand{\kotlincode}[1]{\texttt{\small #1}}
\newcommand{\trythis}[1]{%
  \begin{alertblock}{Try this}
    #1
  \end{alertblock}%
}
\newcommand{\commonmistake}[1]{%
  \begin{exampleblock}{Common mistake}
    #1
  \end{exampleblock}%
}
% Formative in-class checks: question slide, then a dedicated Answer slide
\newcommand{\quizpause}{%
  \par\vspace{0.5em}
  {\small\textit{Pause --- answer (clicker / raise hand / neighbour) before the next slide.}}%
}
\newcommand{\quizbox}[1]{%
  \setbeamercolor{block title}{bg=KeywordBlue,fg=white}%
  \setbeamercolor{block body}{bg=QuizBg,fg=black}%
  \begin{block}{Quiz}
    #1
  \end{block}%
  \setbeamercolor{block title}{bg=LectureNavy,fg=white}%
  \setbeamercolor{block body}{bg=CodeBg,fg=black}%
}
\newcommand{\answerbox}[1]{%
  \setbeamercolor{block title}{bg=AnswerGold,fg=white}%
  \setbeamercolor{block body}{bg=AnswerKeyBg,fg=black}%
  \begin{block}{Answer}
    #1
  \end{block}%
  \setbeamercolor{block title}{bg=LectureNavy,fg=white}%
  \setbeamercolor{block body}{bg=CodeBg,fg=black}%
}

% Empty author/date placeholders for the lecturer to fill
\author{}
\date{}


\usetheme{Madrid}
\usecolortheme{default}

% Clean, professional palette (deep navy + muted teal accent)
\definecolor{LectureNavy}{HTML}{1B3A4B}
\definecolor{LectureAccent}{HTML}{2A9D8F}
\definecolor{LectureGray}{HTML}{4A5568}
\definecolor{CodeBg}{HTML}{F7FAFC}
\definecolor{AlertBg}{HTML}{FFF5F5}
\definecolor{TryBg}{HTML}{F0FFF4}
\definecolor{QuizBg}{HTML}{EBF8FF}
\definecolor{AnswerKeyBg}{HTML}{FFFFF0}
\definecolor{AnswerGold}{HTML}{B7791F}
\definecolor{KeywordBlue}{HTML}{2B6CB0}
\definecolor{StringGreen}{HTML}{276749}
\definecolor{CommentGray}{HTML}{718096}

\setbeamercolor{structure}{fg=LectureNavy}
\setbeamercolor{palette primary}{bg=LectureNavy,fg=white}
\setbeamercolor{palette secondary}{bg=LectureAccent,fg=white}
\setbeamercolor{palette tertiary}{bg=LectureGray,fg=white}
% Madrid puts title/subtitle in a coloured box (inherits palette primary bg).
% Title text must be light on that dark box — not navy-on-navy.
\setbeamercolor{title}{fg=white,bg=LectureNavy}
\setbeamercolor{subtitle}{fg=white,bg=LectureNavy}
\setbeamercolor{author}{fg=LectureNavy}
\setbeamercolor{date}{fg=LectureGray}
\setbeamercolor{institute}{fg=LectureGray}
\setbeamercolor{frametitle}{bg=LectureNavy,fg=white}
\setbeamercolor{block title}{bg=LectureNavy,fg=white}
\setbeamercolor{block body}{bg=CodeBg,fg=black}
\setbeamercolor{block title alerted}{bg=LectureAccent,fg=white}
\setbeamercolor{block body alerted}{bg=TryBg,fg=black}
\setbeamercolor{block title example}{bg=LectureGray,fg=white}
\setbeamercolor{block body example}{bg=AlertBg,fg=black}

\setbeamertemplate{navigation symbols}{}
\setbeamertemplate{itemize items}[circle]
\setbeamertemplate{enumerate items}[default]

\usepackage[T1]{fontenc}
\usepackage[utf8]{inputenc}
\usepackage{lmodern}
\usepackage{booktabs}
\usepackage{tabularx}
\usepackage{graphicx}
\usepackage{hyperref}
\usepackage{listings}
\usepackage{xcolor}
\usepackage{tikz}
\usetikzlibrary{positioning,arrows.meta}

% listings: Kotlin without shell-escape (no built-in Kotlin language)
\lstdefinelanguage{Kotlin}{
  morekeywords={
    abstract,actual,annotation,as,break,by,catch,class,companion,const,
    constructor,continue,crossinline,data,do,dynamic,else,enum,expect,
    external,false,field,file,final,finally,for,fun,get,if,import,in,
    infix,init,inline,inner,interface,internal,is,lateinit,noinline,null,
    object,open,operator,out,override,package,private,protected,public,
    reified,return,sealed,set,super,suspend,tailrec,this,throw,true,try,
    typealias,typeof,val,var,when,where,while
  },
  sensitive=true,
  morecomment=[l]{//},
  morecomment=[s]{/*}{*/},
  morestring=[b]",
  morestring=[b]',
  morestring=[s]{"""}{"""},
}

\lstdefinestyle{kotlinstyle}{
  language=Kotlin,
  basicstyle=\ttfamily\small,
  keywordstyle=\color{KeywordBlue}\bfseries,
  commentstyle=\color{CommentGray}\itshape,
  stringstyle=\color{StringGreen},
  numbers=left,
  numberstyle=\tiny\color{CommentGray},
  numbersep=6pt,
  frame=single,
  rulecolor=\color{LectureGray!40},
  backgroundcolor=\color{CodeBg},
  breaklines=true,
  showstringspaces=false,
  tabsize=2,
  captionpos=b,
  morekeywords={readln,readLine,println,print,listOf,mutableListOf,toInt,toDouble,toLong,toShort,toByte,toChar,toString,split,size},
}
\lstset{style=kotlinstyle}

\newcommand{\kotlincode}[1]{\texttt{\small #1}}
\newcommand{\trythis}[1]{%
  \begin{alertblock}{Try this}
    #1
  \end{alertblock}%
}
\newcommand{\commonmistake}[1]{%
  \begin{exampleblock}{Common mistake}
    #1
  \end{exampleblock}%
}
% Formative in-class checks: question slide, then a dedicated Answer slide
\newcommand{\quizpause}{%
  \par\vspace{0.5em}
  {\small\textit{Pause --- answer (clicker / raise hand / neighbour) before the next slide.}}%
}
\newcommand{\quizbox}[1]{%
  \setbeamercolor{block title}{bg=KeywordBlue,fg=white}%
  \setbeamercolor{block body}{bg=QuizBg,fg=black}%
  \begin{block}{Quiz}
    #1
  \end{block}%
  \setbeamercolor{block title}{bg=LectureNavy,fg=white}%
  \setbeamercolor{block body}{bg=CodeBg,fg=black}%
}
\newcommand{\answerbox}[1]{%
  \setbeamercolor{block title}{bg=AnswerGold,fg=white}%
  \setbeamercolor{block body}{bg=AnswerKeyBg,fg=black}%
  \begin{block}{Answer}
    #1
  \end{block}%
  \setbeamercolor{block title}{bg=LectureNavy,fg=white}%
  \setbeamercolor{block body}{bg=CodeBg,fg=black}%
}

% Empty author/date placeholders for the lecturer to fill
\author{}
\date{}


\usetheme{Madrid}
\usecolortheme{default}

% Clean, professional palette (deep navy + muted teal accent)
\definecolor{LectureNavy}{HTML}{1B3A4B}
\definecolor{LectureAccent}{HTML}{2A9D8F}
\definecolor{LectureGray}{HTML}{4A5568}
\definecolor{CodeBg}{HTML}{F7FAFC}
\definecolor{AlertBg}{HTML}{FFF5F5}
\definecolor{TryBg}{HTML}{F0FFF4}
\definecolor{QuizBg}{HTML}{EBF8FF}
\definecolor{AnswerKeyBg}{HTML}{FFFFF0}
\definecolor{AnswerGold}{HTML}{B7791F}
\definecolor{KeywordBlue}{HTML}{2B6CB0}
\definecolor{StringGreen}{HTML}{276749}
\definecolor{CommentGray}{HTML}{718096}

\setbeamercolor{structure}{fg=LectureNavy}
\setbeamercolor{palette primary}{bg=LectureNavy,fg=white}
\setbeamercolor{palette secondary}{bg=LectureAccent,fg=white}
\setbeamercolor{palette tertiary}{bg=LectureGray,fg=white}
% Madrid puts title/subtitle in a coloured box (inherits palette primary bg).
% Title text must be light on that dark box — not navy-on-navy.
\setbeamercolor{title}{fg=white,bg=LectureNavy}
\setbeamercolor{subtitle}{fg=white,bg=LectureNavy}
\setbeamercolor{author}{fg=LectureNavy}
\setbeamercolor{date}{fg=LectureGray}
\setbeamercolor{institute}{fg=LectureGray}
\setbeamercolor{frametitle}{bg=LectureNavy,fg=white}
\setbeamercolor{block title}{bg=LectureNavy,fg=white}
\setbeamercolor{block body}{bg=CodeBg,fg=black}
\setbeamercolor{block title alerted}{bg=LectureAccent,fg=white}
\setbeamercolor{block body alerted}{bg=TryBg,fg=black}
\setbeamercolor{block title example}{bg=LectureGray,fg=white}
\setbeamercolor{block body example}{bg=AlertBg,fg=black}

\setbeamertemplate{navigation symbols}{}
\setbeamertemplate{itemize items}[circle]
\setbeamertemplate{enumerate items}[default]

\usepackage[T1]{fontenc}
\usepackage[utf8]{inputenc}
\usepackage{lmodern}
\usepackage{booktabs}
\usepackage{tabularx}
\usepackage{graphicx}
\usepackage{hyperref}
\usepackage{listings}
\usepackage{xcolor}
\usepackage{tikz}
\usetikzlibrary{positioning,arrows.meta}

% listings: Kotlin without shell-escape (no built-in Kotlin language)
\lstdefinelanguage{Kotlin}{
  morekeywords={
    abstract,actual,annotation,as,break,by,catch,class,companion,const,
    constructor,continue,crossinline,data,do,dynamic,else,enum,expect,
    external,false,field,file,final,finally,for,fun,get,if,import,in,
    infix,init,inline,inner,interface,internal,is,lateinit,noinline,null,
    object,open,operator,out,override,package,private,protected,public,
    reified,return,sealed,set,super,suspend,tailrec,this,throw,true,try,
    typealias,typeof,val,var,when,where,while
  },
  sensitive=true,
  morecomment=[l]{//},
  morecomment=[s]{/*}{*/},
  morestring=[b]",
  morestring=[b]',
  morestring=[s]{"""}{"""},
}

\lstdefinestyle{kotlinstyle}{
  language=Kotlin,
  basicstyle=\ttfamily\small,
  keywordstyle=\color{KeywordBlue}\bfseries,
  commentstyle=\color{CommentGray}\itshape,
  stringstyle=\color{StringGreen},
  numbers=left,
  numberstyle=\tiny\color{CommentGray},
  numbersep=6pt,
  frame=single,
  rulecolor=\color{LectureGray!40},
  backgroundcolor=\color{CodeBg},
  breaklines=true,
  showstringspaces=false,
  tabsize=2,
  captionpos=b,
  morekeywords={readln,readLine,println,print,listOf,mutableListOf,toInt,toDouble,toLong,toShort,toByte,toChar,toString,split,size},
}
\lstset{style=kotlinstyle}

\newcommand{\kotlincode}[1]{\texttt{\small #1}}
\newcommand{\trythis}[1]{%
  \begin{alertblock}{Try this}
    #1
  \end{alertblock}%
}
\newcommand{\commonmistake}[1]{%
  \begin{exampleblock}{Common mistake}
    #1
  \end{exampleblock}%
}
% Formative in-class checks: question slide, then a dedicated Answer slide
\newcommand{\quizpause}{%
  \par\vspace{0.5em}
  {\small\textit{Pause --- answer (clicker / raise hand / neighbour) before the next slide.}}%
}
\newcommand{\quizbox}[1]{%
  \setbeamercolor{block title}{bg=KeywordBlue,fg=white}%
  \setbeamercolor{block body}{bg=QuizBg,fg=black}%
  \begin{block}{Quiz}
    #1
  \end{block}%
  \setbeamercolor{block title}{bg=LectureNavy,fg=white}%
  \setbeamercolor{block body}{bg=CodeBg,fg=black}%
}
\newcommand{\answerbox}[1]{%
  \setbeamercolor{block title}{bg=AnswerGold,fg=white}%
  \setbeamercolor{block body}{bg=AnswerKeyBg,fg=black}%
  \begin{block}{Answer}
    #1
  \end{block}%
  \setbeamercolor{block title}{bg=LectureNavy,fg=white}%
  \setbeamercolor{block body}{bg=CodeBg,fg=black}%
}

% Empty author/date placeholders for the lecturer to fill
\author{}
\date{}


\title{Lecture 01: First Program and Variables}
\subtitle{Introduction to Programming with Kotlin}

\begin{document}

%----------------------------------------------------------------------
    \begin{frame}
        \titlepage
    \end{frame}

%----------------------------------------------------------------------
    \begin{frame}{Agenda}
        \begin{enumerate}
            \item Your first Kotlin program
            \item Reading from the console
            \item Variables with \kotlincode{val}
            \item Special characters in strings
            \item String templates
        \end{enumerate}
    \end{frame}

%======================================================================
    \section{First program}
%======================================================================

    \begin{frame}{What is Kotlin?}
        \textbf{Kotlin} is a modern, concise programming language that runs on the JVM
        (and elsewhere).
        \vspace{0.8em}
        \begin{itemize}
            \item Designed for clarity and safety (null-safety, sensible defaults)
            \item Fully interoperable with Java
            \item Used for Android, backends, multiplatform apps, and teaching
            \item Official language for Android; popular in industry and education
        \end{itemize}
    \end{frame}

    \begin{frame}[fragile]{Hello, Kotlin}
        The shortest runnable program:
        \vspace{0.4em}
        \begin{lstlisting}
fun main() {
    println("Hello, Kotlin!")
}
        \end{lstlisting}
        \vspace{0.4em}
        \begin{itemize}
            \item \kotlincode{fun main()} --- entry point (no class required for scripts/simple apps)
            \item \kotlincode{println} prints a line and moves to the next line
            \item Save as \kotlincode{Hello.kt}, then compile and run (IDE or \kotlincode{kotlinc})
        \end{itemize}
    \end{frame}

    \begin{frame}{How a Kotlin program runs}
        \begin{block}{Rough pipeline}
            \begin{enumerate}
                \item You write source in a \kotlincode{.kt} file
                \item The Kotlin compiler produces JVM bytecode (or another target)
                \item The runtime executes that bytecode
            \end{enumerate}
        \end{block}
        \vspace{0.6em}
        In this course we focus on writing clear programs, not on build tools.
        Use IntelliJ IDEA / Android Studio / VS~Code with a Kotlin plugin if you like.
    \end{frame}

    \begin{frame}[fragile]{\kotlincode{print} vs \kotlincode{println}}
        \begin{lstlisting}
fun main() {
    print("Hello")
    print(" ")
    println("Kotlin")
    println("Next line")
}
        \end{lstlisting}
        \vspace{0.3em}
        Output:
        \begin{itemize}
            \item \kotlincode{Hello Kotlin} (same line), then
            \item \kotlincode{Next line}
        \end{itemize}
        \vspace{0.4em}
        \trythis{Change the messages and run again. What happens if you use only
        \kotlincode{print}?}
    \end{frame}

    %--- Quiz checkpoint 1 ---
    \begin{frame}{Quiz 1 --- First program}
        \quizbox{What is the entry point of a simple Kotlin application?}
        \begin{enumerate}
            \item \kotlincode{public static void main(String[] args)}
            \item \kotlincode{fun main()}
            \item \kotlincode{class Main}
            \item \kotlincode{System.out.println}
        \end{enumerate}
        \quizpause
    \end{frame}

    \begin{frame}{Answer --- Quiz 1}
        \answerbox{\textbf{2.} \kotlincode{fun main()}}
        \vspace{0.4em}
        Kotlin uses \kotlincode{fun main()} as the program entry point.
        You may also write \kotlincode{fun main(args: Array<String>)}, but for
        beginners the empty-parameter form is enough.
    \end{frame}

%======================================================================
    \section{Console and variables}
%======================================================================

    \begin{frame}{Reading from the console}
        Programs become useful when they react to input.
        \vspace{0.6em}
        \begin{itemize}
            \item \kotlincode{readln()} --- reads one line of text from standard input
            \item Returns a \kotlincode{String} (text)
            \item Prefer \kotlincode{readln()} in modern Kotlin (1.6+); older code may use
              \kotlincode{readLine()}
        \end{itemize}
        \vspace{0.5em}
        \commonmistake{\kotlincode{readLine()} returns \kotlincode{String?} (nullable).
        \kotlincode{readln()} returns a non-null \kotlincode{String} and throws if input ends.}
    \end{frame}

    \begin{frame}[fragile]{Greeting by name}
        \begin{lstlisting}
fun main() {
    print("Enter your name: ")
    val name = readln()
    println("Hello, $name!")
}
        \end{lstlisting}
        \vspace{0.4em}
        \begin{itemize}
            \item Prompt with \kotlincode{print} (stay on the same line)
            \item Store the result in a variable
            \item Use it later in \kotlincode{println}
        \end{itemize}
    \end{frame}

    \begin{frame}{Variables: \kotlincode{val}}
        A \textbf{variable} is a named place that holds a value.
        \vspace{0.6em}
        \begin{itemize}
            \item \kotlincode{val} --- read-only binding (preferred by default)
            \item You give a name, optionally a type, and a value
            \item Names: letters, digits, underscore; start with a letter
            \item Style: \kotlincode{camelCase} (\kotlincode{userName}, \kotlincode{age})
        \end{itemize}
        \vspace{0.5em}
        \begin{block}{Mental model}
            \kotlincode{val name = readln()} means: ``name refers to that string, and
            we will not reassign \kotlincode{name} to something else.''
        \end{block}
    \end{frame}

    \begin{frame}[fragile]{Declaring \kotlincode{val}}
        \begin{lstlisting}
fun main() {
    val greeting = "Welcome"
    val year = 2026
    val pi = 3.14

    println(greeting)
    println(year)
    println(pi)
}
        \end{lstlisting}
        \vspace{0.3em}
        Kotlin \textbf{infers} types when the value is clear:
        \kotlincode{String}, \kotlincode{Int}, \kotlincode{Double}.
        You can write types explicitly: \kotlincode{val year: Int = 2026}.
    \end{frame}

    \begin{frame}[fragile]{Reading numbers (preview)}
        Input is always text. Convert when you need a number:
        \vspace{0.3em}
        \begin{lstlisting}
fun main() {
    print("Enter age: ")
    val age = readln().toInt()
    println("Next year you turn ${age + 1}")
}
        \end{lstlisting}
        \vspace{0.3em}
        \begin{itemize}
            \item \kotlincode{toInt()} parses an integer string
            \item Invalid input (e.g.\ \kotlincode{"abc"}) causes an exception --- we handle that later
        \end{itemize}
        \trythis{Ask for two numbers, convert both with \kotlincode{toInt()}, and print their sum.}
    \end{frame}

    %--- Quiz checkpoint 2 ---
    \begin{frame}{Quiz 2 --- Variables and input}
        \quizbox{Which declaration is correct and idiomatic for a name that will not be reassigned?}
        \begin{enumerate}
            \item \kotlincode{var name = readln()}
            \item \kotlincode{val name = readln()}
            \item \kotlincode{String name = readln()}
            \item \kotlincode{name := readln()}
        \end{enumerate}
        \quizpause
    \end{frame}

    \begin{frame}{Answer --- Quiz 2}
        \answerbox{\textbf{2.} \kotlincode{val name = readln()}}
        \vspace{0.4em}
        Prefer \kotlincode{val} when the binding does not change.
        \kotlincode{var} is for mutable bindings (next lectures).
        Option 3 is Java-style; option 4 is not Kotlin.
    \end{frame}

%======================================================================
    \section{Special characters and templates}
%======================================================================

    \begin{frame}{Special characters in strings}
        Some characters are hard to type ``as themselves'' inside quotes.
        Kotlin uses \textbf{escape sequences} with a backslash:
        \vspace{0.5em}
        \begin{tabular}{ll}
            \toprule
            Escape & Meaning \\
            \midrule
            \kotlincode{\textbackslash n} & new line \\
            \kotlincode{\textbackslash t} & tab \\
            \kotlincode{\textbackslash\textbackslash} & a literal backslash \\
            \kotlincode{\textbackslash"} & a double quote inside a string \\
            \kotlincode{\textbackslash\$} & a literal dollar sign \\
            \bottomrule
        \end{tabular}
    \end{frame}

    \begin{frame}[fragile]{Escapes in action}
        \begin{lstlisting}
fun main() {
    println("Line 1\nLine 2")
    println("Name:\tAda")
    println("She said: \"Hello\"")
    println("Path: C:\\Users\\Student")
}
        \end{lstlisting}
        \vspace{0.3em}
        Raw strings (triple quotes) can hold multiline text with fewer escapes:
        \begin{lstlisting}
val poem = """
First line
Second line
""".trimIndent()
        \end{lstlisting}
    \end{frame}

    \begin{frame}{String templates}
        Build text from variables without awkward concatenation.
        \vspace{0.5em}
        \begin{itemize}
            \item \kotlincode{\$name} --- insert a simple variable
            \item \kotlincode{\$\{expression\}} --- insert any expression
            \item Prefer templates over \kotlincode{+} for readability
        \end{itemize}
    \end{frame}

    \begin{frame}[fragile]{Templates vs concatenation}
        \begin{lstlisting}
fun main() {
    val a = 3
    val b = 5
    // Template (preferred)
    println("$a + $b = ${a + b}")
    // Concatenation (works, noisier)
    println("" + a + " + " + b + " = " + (a + b))
}
        \end{lstlisting}
        \vspace{0.3em}
        Output: \kotlincode{3 + 5 = 8}
        \vspace{0.4em}
        \commonmistake{Writing \kotlincode{\$a + b} only substitutes \kotlincode{a}.
        Use \kotlincode{\$\{a + b\}} for the sum.}
    \end{frame}

    \begin{frame}[fragile]{Putting it together}
        \begin{lstlisting}
fun main() {
    print("First name: ")
    val first = readln()
    print("Year of birth: ")
    val birthYear = readln().toInt()
    val age = 2026 - birthYear
    println("Hello, $first!\nYou turn $age this year.")
}
        \end{lstlisting}
        \trythis{Add a last name and print \kotlincode{Full name: First Last} using a template.}
    \end{frame}

    %--- Quiz checkpoint 3 ---
    \begin{frame}{Quiz 3 --- Templates}
        \quizbox{What does this print? Assume \kotlincode{val n = 4}.\\
        \kotlincode{println("n squared is \$\{n * n\}")}}
        \begin{enumerate}
            \item \kotlincode{n squared is \$\{n * n\}}
            \item \kotlincode{n squared is 4 * 4}
            \item \kotlincode{n squared is 16}
            \item A compile error
        \end{enumerate}
        \quizpause
    \end{frame}

    \begin{frame}{Answer --- Quiz 3}
        \answerbox{\textbf{3.} \kotlincode{n squared is 16}}
        \vspace{0.4em}
        \kotlincode{\$\{n * n\}} evaluates the expression and inserts the result into the string.
    \end{frame}

%======================================================================
    \section{Wrap-up}
%======================================================================

    \begin{frame}{Summary}
        \begin{itemize}
            \item \kotlincode{fun main()} starts the program; \kotlincode{println} writes output
            \item \kotlincode{readln()} reads a line of text from the console
            \item Prefer \kotlincode{val} for bindings that do not change
            \item Escape special characters with \kotlincode{\textbackslash n}, \kotlincode{\textbackslash t}, \ldots
            \item Use string templates: \kotlincode{\$x} and \kotlincode{\$\{...\}}
        \end{itemize}
    \end{frame}

    \begin{frame}{Looking ahead}
        Next lecture: the \kotlincode{Int} type, integer division, mutable variables
        (\kotlincode{var}), floating-point numbers, and other integer sizes
        (\kotlincode{Long}, \kotlincode{Short}, \kotlincode{Byte}).
    \end{frame}

\end{document}
